% !PRINTED 132
\weberpage{132}
%<<CONT

From the continuity of the function $f(x)$ it then follows easily that $f(\xi)=0$; for if $|f(\xi)|>0$, then, since the $\alpha^{(\nu)}$ come arbitrarily close to $\xi$, the absolute values $a^{(\nu)}$ of $f(\alpha^{(\nu)})$ could not sink below every positive value, which is precisely what we have proved for them.


\begin{center}
{\large \S.\ 40.}

\vspace{0.5em}
{\large\gesperrt{Continuity of the Roots.}}
\end{center}

We close this section with the proof of the theorem:

\begin{center}
\gesperrt{The roots of an algebraic equation are continuous functions of the coefficients.}
\end{center}

We have first to explain the meaning of this theorem.

Let
\begin{equation}
 f(x)=x^n+a_1x^{n-1}+a_2x^{n-2}+\cdots
\tag{1}
\end{equation}
be an integral rational function of $x$ of degree $n$. By what has been proved in the preceding paragraphs, $f(x)$ can be decomposed into $n$ linear factors, some of which may be equal to one another. Combining equal factors and denoting the roots by $\alpha,\beta,\gamma,\ldots$, we put
\begin{equation}
 f(x)=(x-\alpha)^a(x-\beta)^b(x-\gamma)^c\cdots,
\tag{2}
\end{equation}
where $a,b,c,\ldots$ are positive whole numbers whose sum is equal to $n$.

The roots $\alpha,\beta,\gamma,\ldots$ will change with the coefficients $a_1,a_2,a_3,\ldots$; the degree of their multiplicity can also become another one.

We denote the changes of $a_1,a_2,
\ldots$ by $\varepsilon_1,\varepsilon_2,\ldots$ and put
\begin{align}
 \varphi(x)&=\varepsilon_1x^{n-1}+\varepsilon_2x^{n-2}+\cdots,
\tag{3}\\
 f(x)+\varphi(x)&=f_1(x).
\tag{4}
\end{align}
We surround the points $\alpha,\beta,\gamma,
\ldots$ by regions of arbitrary smallness, but in such a way that these regions exclude one another; for instance, we enclose the points $\alpha,\beta,\gamma,
\ldots$ by circular peripheries with radii $\varrho,\varrho',\varrho'',\ldots$, and denote these regions by $(\varrho),(\varrho'),(\varrho''),\ldots$.
